\IfFileExists{stacks-project.cls}{%
\documentclass{stacks-project}
}{%
\documentclass{amsart}
}
\usepackage{amssymb}

% For dealing with references we use the comment environment
\usepackage{verbatim}
\usepackage{fancyvrb}
\usepackage{pmboxdraw}
\newenvironment{reference}{\comment}{\endcomment}
%\newenvironment{reference}{}{}
\newenvironment{slogan}{\comment}{\endcomment}
\newenvironment{history}{\comment}{\endcomment}

% For source figures retained as assets
\usepackage{graphicx}

% For commutative diagrams we use Xy-pic
\usepackage[all]{xy}

% We use 2cell for 2-commutative diagrams.
\xyoption{2cell}
\UseAllTwocells

% We use multicol for the list of chapters between chapters
\usepackage{multicol}

% This is generally recommended for better output
\usepackage{lmodern}
\usepackage[T1]{fontenc}
\usepackage{textalpha}

% For cross-file-references
% We prefer the snapshot tmp/xref/#2.aux,
% which latexmkrc replaces atomically, so
% that documents compiled simultaneously
% never read a half-written .aux file.
% The last argument keeps links to #2.pdf.
\usepackage{xr-hyper}
\newcommand{\maybeexternaldocument}[2]{%
  \IfFileExists{#2.aux}{%
    \externaldocument[#1]{#2}%
  }{%
    \PackageWarning{stack}{Missing #2.aux; external references with prefix #1 unavailable}%
  }%
}
\let\externaldocumentorig\externaldocument
\renewcommand{\externaldocument}[2][]{%
  \IfFileExists{tmp/xref/#2.aux}{%
    \externaldocumentorig[#1]{tmp/xref/#2}[#2.pdf]%
  }{%
  \IfFileExists{#2.aux}{%
    \externaldocumentorig[#1]{#2}%
  }{%
    \PackageWarning{stack}{Missing #2.aux; external references with prefix #1 unavailable}%
  }%
  }%
}

% Package for hypertext links:
\usepackage{hyperref}

% For any local file, say "hello.tex" you want to link to please
% use \externaldocument[hello-]{hello}
\externaldocument[introduction-]{introduction}
\externaldocument[conventions-]{conventions}
\externaldocument[sets-]{sets}
\externaldocument[categories-]{categories}
\externaldocument[topology-]{topology}
\externaldocument[sheaves-]{sheaves}
\externaldocument[sites-]{sites}
\externaldocument[stacks-]{stacks}
\externaldocument[fields-]{fields}
\externaldocument[algebra-]{algebra}
\externaldocument[brauer-]{brauer}
\externaldocument[homology-]{homology}
\externaldocument[derived-]{derived}
\externaldocument[simplicial-]{simplicial}
\externaldocument[more-algebra-]{more-algebra}
\externaldocument[smoothing-]{smoothing}
\externaldocument[modules-]{modules}
\externaldocument[sites-modules-]{sites-modules}
\externaldocument[injectives-]{injectives}
\externaldocument[cohomology-]{cohomology}
\externaldocument[sites-cohomology-]{sites-cohomology}
\externaldocument[dga-]{dga}
\externaldocument[dpa-]{dpa}
\externaldocument[sdga-]{sdga}
\externaldocument[hypercovering-]{hypercovering}
\externaldocument[schemes-]{schemes}
\externaldocument[constructions-]{constructions}
\externaldocument[properties-]{properties}
\externaldocument[morphisms-]{morphisms}
\externaldocument[coherent-]{coherent}
\externaldocument[divisors-]{divisors}
\externaldocument[limits-]{limits}
\externaldocument[varieties-]{varieties}
\externaldocument[topologies-]{topologies}
\externaldocument[descent-]{descent}
\externaldocument[perfect-]{perfect}
\externaldocument[more-morphisms-]{more-morphisms}
\externaldocument[flat-]{flat}
\externaldocument[groupoids-]{groupoids}
\externaldocument[more-groupoids-]{more-groupoids}
\externaldocument[etale-]{etale}
\externaldocument[chow-]{chow}
\externaldocument[intersection-]{intersection}
\externaldocument[pic-]{pic}
\externaldocument[weil-]{weil}
\externaldocument[adequate-]{adequate}
\externaldocument[dualizing-]{dualizing}
\externaldocument[duality-]{duality}
\externaldocument[discriminant-]{discriminant}
\externaldocument[derham-]{derham}
\externaldocument[local-cohomology-]{local-cohomology}
\externaldocument[algebraization-]{algebraization}
\externaldocument[curves-]{curves}
\externaldocument[resolve-]{resolve}
\externaldocument[models-]{models}
\externaldocument[functors-]{functors}
\externaldocument[equiv-]{equiv}
\externaldocument[pione-]{pione}
\externaldocument[etale-cohomology-]{etale-cohomology}
\externaldocument[proetale-]{proetale}
\externaldocument[relative-cycles-]{relative-cycles}
\externaldocument[more-etale-]{more-etale}
\externaldocument[trace-]{trace}
\externaldocument[crystalline-]{crystalline}
\externaldocument[spaces-]{spaces}
\externaldocument[spaces-properties-]{spaces-properties}
\externaldocument[spaces-morphisms-]{spaces-morphisms}
\externaldocument[decent-spaces-]{decent-spaces}
\externaldocument[spaces-cohomology-]{spaces-cohomology}
\externaldocument[spaces-limits-]{spaces-limits}
\externaldocument[spaces-divisors-]{spaces-divisors}
\externaldocument[spaces-over-fields-]{spaces-over-fields}
\externaldocument[spaces-topologies-]{spaces-topologies}
\externaldocument[spaces-descent-]{spaces-descent}
\externaldocument[spaces-perfect-]{spaces-perfect}
\externaldocument[spaces-more-morphisms-]{spaces-more-morphisms}
\externaldocument[spaces-flat-]{spaces-flat}
\externaldocument[spaces-groupoids-]{spaces-groupoids}
\externaldocument[spaces-more-groupoids-]{spaces-more-groupoids}
\externaldocument[bootstrap-]{bootstrap}
\externaldocument[spaces-pushouts-]{spaces-pushouts}
\externaldocument[spaces-chow-]{spaces-chow}
\externaldocument[groupoids-quotients-]{groupoids-quotients}
\externaldocument[spaces-more-cohomology-]{spaces-more-cohomology}
\externaldocument[spaces-simplicial-]{spaces-simplicial}
\externaldocument[spaces-duality-]{spaces-duality}
\externaldocument[formal-spaces-]{formal-spaces}
\externaldocument[restricted-]{restricted}
\externaldocument[spaces-resolve-]{spaces-resolve}
\externaldocument[formal-defos-]{formal-defos}
\externaldocument[defos-]{defos}
\externaldocument[cotangent-]{cotangent}
\externaldocument[examples-defos-]{examples-defos}
\externaldocument[algebraic-]{algebraic}
\externaldocument[examples-stacks-]{examples-stacks}
\externaldocument[stacks-sheaves-]{stacks-sheaves}
\externaldocument[criteria-]{criteria}
\externaldocument[artin-]{artin}
\externaldocument[quot-]{quot}
\externaldocument[stacks-properties-]{stacks-properties}
\externaldocument[stacks-morphisms-]{stacks-morphisms}
\externaldocument[stacks-limits-]{stacks-limits}
\externaldocument[stacks-cohomology-]{stacks-cohomology}
\externaldocument[stacks-perfect-]{stacks-perfect}
\externaldocument[stacks-introduction-]{stacks-introduction}
\externaldocument[stacks-more-morphisms-]{stacks-more-morphisms}
\externaldocument[stacks-geometry-]{stacks-geometry}
\externaldocument[moduli-]{moduli}
\externaldocument[moduli-curves-]{moduli-curves}
\externaldocument[examples-]{examples}
\externaldocument[exercises-]{exercises}
\externaldocument[guide-]{guide}
\externaldocument[desirables-]{desirables}
\externaldocument[coding-]{coding}
\externaldocument[obsolete-]{obsolete}
\externaldocument[fdl-]{fdl}
\externaldocument[index-]{index}

% Heap Project documents
\externaldocument[k3-]{k3}
\externaldocument[torelli-]{torelli}
\externaldocument[linear-algebra-]{linear-algebra}
\externaldocument[groups-]{groups}
\externaldocument[complex-analysis-]{complex-analysis}
\externaldocument[functional-analysis-]{functional-analysis}
\externaldocument[categories-]{categories}
\externaldocument[commutative-algebra-]{commutative-algebra}
\externaldocument[ringed-spaces-]{ringed-spaces}
\externaldocument[homological-algebra-]{homological-algebra}
\externaldocument[lie-algebras-]{lie-algebras}
\externaldocument[representation-theory-]{representation-theory}
\externaldocument[smooth-manifolds-]{smooth-manifolds}
\externaldocument[differential-topology-]{differential-topology}
\externaldocument[algebraic-topology-]{algebraic-topology}
\externaldocument[differential-geometry-]{differential-geometry}
\externaldocument[algebraic-geometry-]{algebraic-geometry}
\externaldocument[symplectic-geometry-]{symplectic-geometry}
\externaldocument[complex-geometry-]{complex-geometry}
\externaldocument[hodge-theory-]{hodge-theory}
\externaldocument[%
variations-of-hodge-structures-]{%
variations-of-hodge-structures}
\externaldocument[%
homological-mirror-symmetry-]{%
homological-mirror-symmetry}
\maybeexternaldocument{deformations-}{deformations}
\externaldocument[vertex-algebras-]{vertex-algebras}
\externaldocument[springer-]{springer}
\externaldocument[infty-categories-]{infty-categories}
\externaldocument[seminars-differential-geometry-]{%
seminars-differential-geometry}
\externaldocument[seminars-complex-geometry-]{%
seminars-complex-geometry}
\externaldocument[seminars-algebraic-geometry-]{%
seminars-algebraic-geometry}
\externaldocument[seminars-analysis-pde-]{%
seminars-analysis-pde}
\externaldocument[surfaces-]{surfaces}
\externaldocument[birational-maps-commutative-algebra-]{birational-maps-commutative-algebra}
\externaldocument[cremona-transformations-]{cremona-transformations}
\externaldocument[derived-categories-of-sheaves-]{derived-categories-of-sheaves}
\externaldocument[geometric-invariant-theory-]{geometric-invariant-theory}
\externaldocument[geometric-stability-conditions-and-group-actions-]{geometric-stability-conditions-and-group-actions}
\externaldocument[moduli-spaces-of-sheaves-]{moduli-spaces-of-sheaves}
\externaldocument[bridgeland-stability-general-theory-]{bridgeland-stability-general-theory}
\externaldocument[hilbert-schemes-]{hilbert-schemes}
\externaldocument[noncommutative-abelian-surfaces-and-kummer-type-hyperkahler-varieties-]{noncommutative-abelian-surfaces-and-kummer-type-hyperkahler-varieties}
\externaldocument[mumford-tate-groups-in-hodge-theory-]{mumford-tate-groups-in-hodge-theory}
\externaldocument[hodge-birational-atoms-2026-]{hodge-birational-atoms-2026}
\externaldocument[xxii-egd-2026-]{%
xxii-egd-2026}
\externaldocument[pde-]{pde}
\externaldocument[probability-]{probability}

\externaldocument[linguistics-]{linguistics}

\externaldocument[%
16th-alga-meeting-2026-]{%
16th-alga-meeting-2026}

\externaldocument[%
iv-jornadas-lefschetz-2026-]{%
iv-jornadas-lefschetz-2026}

\externaldocument[reading-the-masters-]{%
reading-the-masters}

\externaldocument[real-analysis-]{%
real-analysis}

\externaldocument[optimal-transport-]{%
optimal-transport}

\externaldocument[history-]{history}

\externaldocument[seiberg-witten-]{seiberg-witten}

\externaldocument[history-philosophy-biology-]{history-philosophy-biology}

\externaldocument[classical-mechanics-]{classical-mechanics}
\externaldocument[electromagnetism-]{electromagnetism}
\externaldocument[waves-and-optics-]{waves-and-optics}
\externaldocument[wif-]{waves-in-fluids-iv}
\externaldocument[thermal-physics-]{thermal-physics}
\externaldocument[quantum-mechanics-]{quantum-mechanics}
\externaldocument[special-relativity-]{special-relativity}
\externaldocument[physics-]{physics}

\externaldocument[chemistry-]{chemistry}

\externaldocument[molecular-cell-biology-]{molecular-cell-biology}

\externaldocument[microbiology-]{microbiology}

\externaldocument[protistology-]{protistology}

\externaldocument[mycology-]{mycology}

\externaldocument[botany-]{botany}

\externaldocument[zoology-]{zoology}

\externaldocument[animal-histology-]{animal-histology}

\externaldocument[animal-physiology-]{animal-physiology}

\externaldocument[systematics-biogeography-]{systematics-biogeography}

\externaldocument[ecology-natural-resources-]{ecology-natural-resources}

\externaldocument[earth-sciences-]{earth-sciences}

\externaldocument[biostatistics-]{biostatistics}

\externaldocument[mathematical-systems-biology-]{mathematical-systems-biology}

\externaldocument[cimpa-el-salvador-2024-]{cimpa-el-salvador-2024}

\externaldocument[anthropology-]{anthropology}

\externaldocument[yang-mills-theory-]{yang-mills-theory}

\externaldocument[seminars-algebra-]{seminars-algebra}

\externaldocument[seminars-physics-]{seminars-physics}

\externaldocument[seminars-computing-]{seminars-computing}

\externaldocument[oktobermat-2026-]{oktobermat-2026}

\externaldocument[saet-2026-]{saet-2026}

\externaldocument[seminars-symplectic-geometry-]{seminars-symplectic-geometry}

\externaldocument[seminars-dynamics-]{seminars-dynamics}

\externaldocument[oktobermat-2024-]{oktobermat-2024}

\externaldocument[quantum-groups-]{quantum-groups}

\externaldocument[legal-2027-]{legal-2027}

\externaldocument[derived-categories-]{derived-categories}

% Theorem environments.
%
\theoremstyle{plain}
\newtheorem{theorem}[subsection]{Theorem}
\newtheorem{proposition}[subsection]{Proposition}
\newtheorem{lemma}[subsection]{Lemma}
\newtheorem{corollary}[subsection]{Corollary}

\theoremstyle{definition}
\newtheorem{definition}[subsection]{Definition}
\newtheorem{example}[subsection]{Example}
\newtheorem{exercise}[subsection]{Exercise}
\newtheorem{situation}[subsection]{Situation}

\theoremstyle{remark}
\newtheorem{remark}[subsection]{Remark}
\newtheorem{remarks}[subsection]{Remarks}

\numberwithin{equation}{subsection}

% Macros
%
\def\lim{\mathop{\mathrm{lim}}\nolimits}
\def\colim{\mathop{\mathrm{colim}}\nolimits}
\def\Spec{\mathop{\mathrm{Spec}}}
\def\Hom{\mathop{\mathrm{Hom}}\nolimits}
\def\Ext{\mathop{\mathrm{Ext}}\nolimits}
\def\SheafHom{\mathop{\mathcal{H}\!\mathit{om}}\nolimits}
\def\SheafExt{\mathop{\mathcal{E}\!\mathit{xt}}\nolimits}
\def\Sch{\mathit{Sch}}
\def\Mor{\mathop{\mathrm{Mor}}\nolimits}
\def\Ob{\mathop{\mathrm{Ob}}\nolimits}
\def\Sh{\mathop{\mathit{Sh}}\nolimits}
\def\NL{\mathop{N\!L}\nolimits}
\def\CH{\mathop{\mathrm{CH}}\nolimits}
\def\proetale{{pro\text{-}\acute{e}tale}}
\def\etale{{\acute{e}tale}}
\def\QCoh{\mathit{QCoh}}
\def\Ker{\mathop{\mathrm{Ker}}}
\def\Im{\mathop{\mathrm{Im}}}
\def\Coker{\mathop{\mathrm{Coker}}}
\def\Coim{\mathop{\mathrm{Coim}}}

% Boxtimes
%
\DeclareMathSymbol{\boxtimes}{\mathbin}{AMSa}{"02}

%
% Macros for moduli stacks/spaces
%
\def\QCohstack{\mathcal{QC}\!\mathit{oh}}
\def\Cohstack{\mathcal{C}\!\mathit{oh}}
\def\Spacesstack{\mathcal{S}\!\mathit{paces}}
\def\Quotfunctor{\mathrm{Quot}}
\def\Hilbfunctor{\mathrm{Hilb}}
\def\Curvesstack{\mathcal{C}\!\mathit{urves}}
\def\Polarizedstack{\mathcal{P}\!\mathit{olarized}}
\def\Complexesstack{\mathcal{C}\!\mathit{omplexes}}
% \Pic is the operator that assigns to X its picard group, usage \Pic(X)
% \Picardstack_{X/B} denotes the Picard stack of X over B
% \Picardfunctor_{X/B} denotes the Picard functor of X over B
\def\Pic{\mathop{\mathrm{Pic}}\nolimits}
\def\Picardstack{\mathcal{P}\!\mathit{ic}}
\def\Picardfunctor{\mathrm{Pic}}
\def\Deformationcategory{\mathcal{D}\!\mathit{ef}}


\begin{document}

\title{Seminars in Physics}
\maketitle

\phantomsection
\label{section-phantom}

\tableofcontents

\section{Neutrinos}
\label{section-neutrinos}
%% speaker: Hiroshi Nunokawas, PUC-Rio
%% place: Friday Seminar
%% date: June 27, 2025

\noindent
Hiroshi Nunokawas, PUC-Rio.
Friday Seminar, Seminar Name.
June 27, 2025.

\medskip\noindent
{\bf Abstract.} Hiroshi will come and tell us
everything we (not) wanted to know
about these mysterious particles, and are not
going to be afraid to ask. In
particular, about the neutrino oscillation,
and the great matrices.

\bigskip\noindent

Protons and neutrons have very similar mass
of $m_p\approx 940$ MeV, while
electrons have mass of $m_e \approx 0.5$ MeV.
MeV is $10^{6}$ electronvolts,
where one eV is approximately $1.6\times
10^{-19}$ J. This is standard in high
energy physics, they use electronvolts
instead of Joules. Recall that
2J=1N$\times$1m.

Most of the things we see are protons since
they are so much larger than
electrons. But protons nor neutrons are
elementary particles.

Here's the standard model:

\medskip

\begin{tabular}{c c c c}
Quarks & $\begin{bmatrix} u\\d
\end{bmatrix}_L $ & $\begin{bmatrix} c\\s
\end{bmatrix}_L$ & $\begin{bmatrix} t\\b
\end{bmatrix}_L$\\
Leptons & $\begin{bmatrix} \nu_e\\e^-
\end{bmatrix}_L $ & $\begin{bmatrix}
\nu_\mu\\ \mu^-
\end{bmatrix}_L$ & $\begin{bmatrix}
\nu_\tau\\ \tau^- \end{bmatrix}_L$\\
Generation & 1st & 2nd & 3rd\\
Bosons & $g$, $\gamma$ , $\omega^\pm$, $z$ \\
Higgs Bosson & $H$
\end{tabular}

\medskip

It is very particular that nature repeats
itself three times. The $L$ in those
matrix actually means left-handed, and
accounts for chirality. Only left-handed
fermions have weak interaction. Right-handed
have electromagnetic interaction,
gravitational interaction, but not weak
interaction.

And then there's neutrinos. They have
negative helicity (chirality). Being
left-handed, mathematically, means to have
helicity $-1$. I think this means
that the spin is left-handed. But chirality
and helicity are not the same:
helicity is observer-dependent, and chirality
is not. Almost all neutrinos we
can see (\% 99.99999…) have negative
helicity, but not all of them.

Consider the following:
$$
n+\nu_e\leftrightarrow p+e^-
$$
But it's not completely correct: we'd better
put $d$ instead of $n$, and $u$
instead of $p$: the $d$ and $u$ quarks,
instead of the neutrons and protons.

Now consider the following reaction: a
neutron decays into a proton, an electron
and an antineutrino:
$$
n\to p+e^+\overline{\nu}_e
$$
Protons is very stable, that's why we are
here. But neutron decays in only 15
minutes.

By experimental data, we can conclude that
neutrinos' mass is consistent with
zero. But if they have mass, it should be
much smaller than the electron's $m_v
\leq 0.5$ eV. And the electron is already the
lightest fermion!

If the mass of the neutrino was zero, i.e.
$m_\nu=0$, then $v_\nu=c$ in vacuum,
which would imply that
$$
\xymatrix{
\nu_e\ar[r]^{L}&\nu_e\ar[r]&\nu_e\\
0:00&0:00&0:00
}
$$
meaning: time doesn't pass! And this means
the state of the particle cannot
change.

\section{Second Quantization}
\label{section-erdnigor-second-quantization}
%% speaker: Altan Erdnigor, IMPA
%% place: Seminário das Sextas, PUC-Rio
%% date: October 18, 2024

\noindent
Altan Erdnigor, IMPA.
Semin\'ario das Sextas, PUC-Rio.
October 18, 2024.

\medskip\noindent
{\bf Abstract.}
The talk will cover a survey by Yuri Neretin
of a famous eponymous book by Felix Berezin.
Neretin's overview is short and clean, with
no proofs (they are all exercises in
functional analysis), providing an abundance
of references.

The topics are: finite-dimensional and
infinite-dimensional bosonic Fock space
(\href{https://en.wikipedia.org/wiki/Segal%E2%80%93Bargmann_space}{Segal-Bargmann
space}
(1961, 1963)) and fermionic Fock spaces, and
the Weil representation of the orthogonal and
symplectic groups.

The main theorem of the paper is that the
representations do exist. Mostly it's about
giving the right definitions. If we don't
discuss the proofs, I can fit it under one
hour.

\bigskip\noindent
{\bf Plan.}
\begin{enumerate}
\item Bossonic Fock space $F_n\cong
L^2(\mathbb{R}^n)$. What is it, how to work
with it. This will be the quantization of
ordinary space $\mathbb{R}^{2}$.

\item Metaplectic representation (Weil
representation) $\mathsf{Sp}(\mathbb{R})$
(symplectic matrices). The projective action
$\mathsf{Sp}_{2n}(\mathbb{R})\curvearrowright
F_n$. You will see how strange it is that
this representation ?

\item Fermionic Fock space $\Lambda_n\cong
\Lambda^\bullet(\mathbb{C}^n)$.

\item Spinor represenation. Fermionic Fock
space also admits a projective representation
called spinor represenation: $\mathsf{SO}(2n)
\curvearrowright \Lambda_n$.

\item $n\to \infty$. Gives rise to $F_\infty,
\Lambda_\infty$.
\end{enumerate}

Big formulas give rise to representations.

\medskip\noindent{\bf Fock space.}
Consider $\mathbb{C}^{n}$. define
$$F_n=\left\{\text{$f(z)$-entire functions
$\mathbb{C}^{n}\to \mathbb{C}$}\Big|
\int_{\mathbb{C}^{n}}(f(x))^2e^{-|z|^2}
d\lambda(z)<\infty\right\}$$

We have an inner product:
$$\left\langle{} f,g \right\rangle{}
:=\int_{\mathbb{C}^{n}}f(z)\overline{g(z)}
e^{|z|^2}d
\lambda(z)$$

\begin{proposition}
\label{proposition-erdnigor-monomials}
$z^k:=z_1^{k_1}\cdot\ldots\cdot z_n^{k_n}$,
where $k=(k_1,\ldots,k_n)$, form an
orthogonal basis (of monomials) for $F_n$.
Thus $F_n$ is a Hilbert space. Moreover,
$\|z^k\|^2=k_1!\ldots k_n!$.
\end{proposition}

\medskip\noindent{\bf Question.}
Why is this basis orthogonal--- why do two
different monomials integrate to zero?

\begin{remark}
\label{remark-erdnigor-folland}
See {\it Harmonica analysis in phase space}
by Folland (1988).
\end{remark}

\begin{definition}
\label{definition-erdnigor-bt}
$b_T(z):=\operatorname{exp}(zTz^+)$ for $T$
symmetric.
\end{definition}

\medskip\noindent{\bf Claim.}
$b_T\in F_n \iff \|T\|<1$.

\begin{exercise}
\label{exercise-erdnigor-bt}
$\left\langle{} b_T,b_S \right\rangle{} =\det
( (1-TS^*)^{-\frac{1}{2}} )$ {\bf Hint:}
Appendix 1 in Folland. This implies that
$$\|b_T\|=\det (1-T T^*
)^{-\frac{1}{4}}=\sqrt{\det (1-T
T^*)^{-\frac{1}{2}}}$$
\end{exercise}

\medskip\noindent
{\bf Altan's talk exercise 2: Fock space is
reproducing kernel space.}

\begin{definition}[Coherent states]
\label{definition-erdnigor-coherent}
For $a\in \mathbb{C}^{n}$,
$$\varphi_a(z):=\operatorname{exp}
(z_1\overline{a_1}+\ldots
+z_n\overline{a_n})$$
\end{definition}

\medskip\noindent{\bf Claim.}
For $f\in F_n$,
$$\left\langle{} f,\varphi_a \right\rangle{}
=f(a)$$

\begin{remark}
\label{remark-erdnigor-rkhs}
This says that for every $a\in\mathbb{C}^{n}$
the functional $F_n\ni f\mapsto
f(z)\in\mathbb{C}$ is bounded since it is
represented by $\varphi_a$ (Riesz
representation theorem). This is called
{\it\bf Reproducing Kernel Hilbert Space}.
\end{remark}

\begin{proof}[Proof (Exercise)]
Let's consider the case $n=1$. Then
$\varphi_a=\operatorname{exp}(z\bar{a})$ and
$$\left\langle{} f,\varphi_a \right\rangle{}
=\int_{\mathbb{C}}f(z)
\overline{\operatorname{exp}(za)
}e^{-|z|^2}d \lambda(z)$$
But I'm not sure how to compute this
integral. Upon difficulties let's suppose
$f\equiv1$. Perhaps split complex exponential
into real and imaginary parts? Then we get
real-valued integrals over
$\mathbb{R}^{2}=\mathbb{C}$:
$$
\begin{aligned}
\int_{\mathbb{C}}\overline{\operatorname{exp}
(za)}
e^{-|z|^2}d
\lambda(z)&=\int_{\mathbb{R}^{2}}
\overline{e^{\operatorname{Re}(za)}\big(\cos
\operatorname{Im}(za)+i
\operatorname{Im}(za)\big)}e^{-|z|^2}d
\lambda(z)\\
&=
\int_{\mathbb{R}^{2}}
\overline{e^{\operatorname{Re}(za)}\cos
\operatorname{Im}(za)}e^{-|z|^2}d
\lambda(z)+\int_{\mathbb{R}^{2}}
\overline{e^{\operatorname{Re}(za)}i
\operatorname{Im}(za)}e^{-|z|^2}d \lambda(z)
\end{aligned}
$$
Not sure$\ldots$

In the general case we shall have
$$
\begin{aligned}
\left\langle{} f,\varphi_a \right\rangle{}
&=\int_{\mathbb{C}^n}f(z)
\overline{\varphi_a(z)}e^{-|z|^2}d\lambda(z)
\\
&=\int_{\mathbb{C}^{n}}f(z)
\overline{\operatorname{exp}(z_1\overline{a}
_1+\ldots
+z_n\overline{a}_n)}e^{-|z|^2}d \lambda(z)
\end{aligned}
$$
Maybe take this somehow to a line integral
around $a$ to Cauchy integral
formula$\ldots$?
\end{proof}

\begin{remark}
\label{remark-erdnigor-evaluation}
Evaluation at a point is a bounded functional
(this is not the case in general  $L^2$ space
because you can have functions with
arbitrarily large values).
\end{remark}

\medskip\noindent{\bf Question.}
What are the operators $A\curvearrowright
F_n$?

\medskip\noindent{\bf Answer.}
$(\operatorname{Af})z=\int_{\mathbb{C}^{n}}
k(z,\bar{u})f(u)e^{-|u|^2}
d \lambda(u)$, where $K(z,u)$ is holomorphic
in  $z$ and antiholomorphic in $u$. The
integral converges absolutely for all $f \in
F_n$.

\begin{proof}
We need to compute the kernel. So define
$$c_ke=\left\langle{} Az^{\mathbf{k}},z^\ell
\right\rangle{}$$
then
$$k(z,\bar{u}
)=\sum_{k,\ell}c_{k\ell}\frac{z^k}{k!}
\frac{\bar{u}
^\ell}{\ell!}$$
which means that the kernel is symmetric. So
this means that restricting the kernel to
diagonal is not holomorphic anymore, kind of
real analytic (?).

Alternatively, $k(a,b)=\left\langle{}
A\varphi_b,\varphi_a \right\rangle{}$.
\end{proof}

\medskip\noindent{\bf Question (Dani).}
Why do we need to compute the kernel? We want
to show that the operator $A\curvearrowright
F_n$ produces a (projective) action, right?

\medskip\noindent{\bf OK.}
so we defined Fock space and the operators
that act nicely on it.

\medskip\noindent{\bf Metaplectic
representation.}

\begin{definition}
\label{definition-erdnigor-sp}
$$\mathsf{Sp}(2n,\mathbb{R}:=\left(
h\in\operatorname{Mat}_{2n\times 2n}:h
\begin{pmatrix} 0&1\\-1&0
\end{pmatrix}h^+=\begin{pmatrix} 0&1\\-1&0
\end{pmatrix}   \right)$$
\end{definition}

Let
$$J=\frac{1}{\sqrt{2} }\begin{pmatrix}
1&i\\i&1 \end{pmatrix} , \qquad g:=Jg
J^{-1}$$
$$g=\begin{pmatrix} Q&\psi\\\bar{\psi}
&\bar{\phi}  \end{pmatrix} ,\qquad g
\begin{pmatrix} 1&0\\0&-1 \end{pmatrix} g^*
=\begin{pmatrix} 1&0\\0&-1 \end{pmatrix}$$

Also
$$\mathsf{Sp}(2n,\mathbb{R})\cong
\mathsf{U}(n,n) \cap
\mathsf{Sp}(2n,\mathbb{C})$$
Now
$$\left(W\begin{pmatrix}
\phi&\psi\\\bar{\psi} &\bar{\phi}
\end{pmatrix}
f\right)(z)=\int_{\mathbb{C}^{n}}
\operatorname{exp}
\left\{ \frac{1}{2}(z,\bar{u}
)\begin{pmatrix} \bar{\psi}
\phi^{-1}&(\phi^t)^{-1}\\\phi^{-1}&-\phi^{-1}
\psi
\end{pmatrix} \begin{pmatrix} z^t\\\bar{u} ^t
\end{pmatrix}  \right\} f(u)e^{-|u|^2}d
\lambda(u).$$

Who is this guy? It's a {\it\bf Berezin's
integral}. So that is the definition of $W$,
which is the representation of the
mataplectic group
$W:\mathsf{Sp}_{2n}(\mathbb{R})
\curvearrowright
F_n$ we wanted.

Now let's explain what is a projective
representation.

\begin{theorem}
\label{theorem-erdnigor-projective}
\begin{enumerate}

\item
$W(\cdot )$ are unitary up to scalar. More
precisely,
$$\det (\phi^*\phi)^{-\frac{1}{4}} W
\begin{pmatrix} \phi&\psi\\\bar{\psi}
&\bar{\phi}  \end{pmatrix}$$
are unitary.

\item $W(\cdot )$ define a projective
representation of
$\mathsf{Sp}_{2n}(\mathbb{R})\curvearrowright
F_n$.

Now let's finish with this nice formula for
the cocycle (how to multiply these
operators?:

$$W \begin{pmatrix} \phi_1
&\psi_1\\\bar{\psi}_1 &\bar{\phi}_1
\end{pmatrix}W \begin{pmatrix}
\phi_2&\psi_2\\\bar{\psi}_2&\bar{\phi}_2
\end{pmatrix} =\det
(1+\phi_1^{-1}\psi_1\bar{\psi}_2\phi_2^{-1})
^{-\frac{1}{2}}
W\left( \begin{pmatrix}
\phi_1&\psi_1\\\bar{\psi}_1&\bar{\phi}_1
\end{pmatrix}  \right)$$
\end{enumerate}
\end{theorem}

\medskip\noindent{\bf Question (Dani).}
So what is projective representation?

\medskip\noindent
{\bf Editorial note.}
The notes are unfinished and end with this
question. Colour marks of the original
(for questions and doubts) were dropped,
the section titles Plan, Fock space and
Metaplectic representation became
bold headings, the old claims and
questions became bold run-in headings,
and the angle brackets of the inner
products were rewritten as required by
HP style. Typos and doubtful formulas,
including the last formula of the
theorem, are as in the original.

\section{Feynman checkers: quantum theory on
a checkered paper}
\label{section-skopenkov-feynman-checkers}
%% speaker: Mikhail Skopenkov, King Abdullah
%% speaker: University of Science and
%% speaker: Technology
%% place: Seminário das Sextas
%% date: February 27, 2025

\noindent
Mikhail Skopenkov, King Abdullah University
of Science and Technology.
Semin\'ario das Sextas.
February 27, 2025.

\medskip\noindent
{\bf Abstract.}
We survey the most elementary model of
electron motion introduced by Feynman. In
this game, a checker moves on a checkerboard
by simple rules, and we count the turns.
Feynman checkers are also known as a
one-dimensional quantum walk, or an Ising
model at an imaginary temperature, or a
six-vertex model with complex weights. We
mathematically solve a Feynman problem from
1965, which was to prove that the discrete
model (for large time, small average
velocity, and small lattice step) is
consistent with the continuum one. We get
further asymptotic properties of the model
for large time, improving the results by
Sunada and Tate from 2012 towards the
ultimate goal of a uniform asymptotic formula
in the whole range.

\medskip\noindent
{\bf Editorial note.}
The old notes contain only the title,
speaker and abstract of this talk, no notes.
The old header gives the series
(Seminário das Sextas) but not the
institution.

\bibliography{my}
\bibliographystyle{amsalpha}

\end{document}
